\documentclass{beamer}

\usetheme{default}

% Title page
\title{Cubesphere overview}

\begin{document}

\begin{frame}
    \titlepage
\end{frame}

\begin{frame}{Cube panels}
\begin{figure}
    \includegraphics[scale=0.5]{img/panel_numbering}
    \caption{Panel numbering: This shows how the panels (cube faces) are numbered and connected to each other.}
\end{figure}
\end{frame}


\begin{frame}{Panel orientation}
\begin{itemize}
    \item The orientation is local to the tile, not geographical
    \item Note that the North of tile~2 is connected to the North of tile~4
\end{itemize}
\begin{figure}
    \includegraphics[scale=0.9]{img/tile_orientation}
\end{figure}
\end{frame}

\begin{frame}{Partitioning for MPI}
\begin{itemize}
    \item The tiles all have identical sizes, horizontally and vertically.
    \item They are grouped by panel.
\end{itemize}
\begin{figure}
    \includegraphics[scale=0.6]{img/multi_tile_numbering}
    \caption{Tile numbering: How tiles are assigned to MPI ranks.}
\end{figure}
\end{frame}

\begin{frame}{Data storage (single variable)}
\begin{itemize}
    \item Data is stored by row \emph{of points}, not elements
\end{itemize}
\begin{figure}
    \includegraphics[scale=0.8]{img/variable_storage}
\end{figure}
\end{frame}

\end{document}
